% from ega3-7-fr.tex
\oldpage[III]{221}
séparé, $\Gamma_f$ est une immersion fermée (I, 5.4.3) donc quasi-compacte
(I, 6.6.4). Si au contraire $X$ est localement noethérien, comme $Y$ est affine
et $f$ quasi-compact, l'espace sous-jacent à $X$ est quasi-compact, donc
noethérien, et on peut de nouveau appliquer (I, 6.6.4) pour prouver que
$\Gamma_f$ est quasi-compacte.

\noindent\textbf{$(\mathbf{Err}_{\mathrm{III}},\,19)$} L'énoncé de (II,
5.2.2) est incorrect, les conditions $b)$, $c)$ et $c')$ n'étant pas locales
sur $Y$ lorsqu'on ne sait pas si pour un ouvert $U$ de $Y$, un
$(\mathscr O_X|f^{-1}(U))$-Module quasi-cohérent est restriction à $f^{-1}(U)$
d'un $\mathscr O_X$-Module quasi-cohérent. Il faut donc remplacer dans la ligne
16 de la p. 98 « \emph{Soit $f:X\to Y$ un morphisme séparé quasi-compact} »
par : « \emph{Soient $X$, $Y$ deux préschémas tels que $X$ soit un schéma ou
que l'espace sous-jacent à $X$ soit localement noethérien, et soit
$f:X\to Y$ un morphisme quasi-compact.} » Dans la démonstration, on observera
que les hypothèses entraînent que $f$ est \emph{séparé} lorsqu'on suppose que
$X$ est un schéma, en vertu de (I, 5.5.5 et 5.5.8) ; on peut donc appliquer
(I, 9.2.2) dans les deux cas.

\noindent\textbf{$(\mathbf{Err}_{\mathrm{III}},\,20)$} Dans (II, 6.2.3),
ajouter, après la ligne 18 du bas de la p. 115 :

\emph{On dit qu'un morphisme $f:X\to Y$ est quasi-fini en un point $x\in X$
s'il existe un voisinage ouvert affine $V$ de $y=f(x)$ et un voisinage ouvert
affine $U$ de $x$ tels que $f(U)\subset V$ et que le morphisme $U\to V$
restriction de $f$ soit quasi-fini. On dit qu'un morphisme $f:X\to Y$ est
localement quasi-fini s'il est quasi-fini en tout point de $X$.}

\noindent\textbf{$(\mathbf{Err}_{\mathrm{III}},\,21)$} Dans (II, 6.4.3), la
démonstration est incorrecte, les éléments d'un sous-$A$-module de type fini de
$K$ n'étant pas nécessairement entiers sur $A$. Remplacer les 9 dernières lignes
de la p. 121 par le texte suivant : Comme on peut supposer que $E$ est sans
torsion, donc fidèle, $E$ est aussi un $A[u]$-module fidèle ($A$ étant plongé
dans l'anneau des endomorphismes de $E$). Comme $E$ est un $A$-module de type
fini, il en résulte que $u$ est \emph{entier} sur $A$ (Bourbaki,
\emph{Alg. comm.}, chap. V, \S~1, n\textsuperscript{o}~1, lemme 1). On en
conclut immédiatement que les valeurs propres de $u\otimes1$ (dans une clôture
algébrique de $K$) sont des éléments \emph{entiers} sur $A$, et il en est donc
de même des $\sigma_i(u)$.

\noindent\textbf{$(\mathbf{Err}_{\mathrm{III}},\,22)$} Dans
$(0_{\mathrm{III}},\,9.1.1)$, ligne 18 du bas de la p. 12, supprimer
« ouverte ».

\noindent\textbf{$(\mathbf{Err}_{\mathrm{III}},\,23)$} Dans
$(0_{\mathrm{III}},\,11.7.3)$, ligne 10 de la p. 43, après $\mathsf C''$,
ajouter : tel que pour tout objet projectif $P$ de $\mathsf C'$ (resp. tout
objet projectif $P'$ de $\mathsf C''$) le foncteur
$A'\rightsquigarrow T(P,A')$ (resp. $A\rightsquigarrow T(A,P')$) soit exact
dans $\mathsf C'$ (resp. $\mathsf C$).

\noindent\textbf{$(\mathbf{Err}_{\mathrm{III}},\,24)$} L'énoncé de
$(0_{\mathrm{III}},\,13.7.7)$ est inexact, et doit être modifié comme suit :
ligne 6 de la p. 78, supprimer « et
$(\mathrm R^{n+1}T(A_k))_{k\in\mathbf Z}$ » ; ligne 7 de la p. 78, remplacer
« $\mathrm R'^nT(A)$ est un $S$-module de type fini » par
« $\mathrm R'^nT(A)$ et $\mathrm R'^{n+1}T(A)$ sont des $S$-modules de type
fini » ; lignes 12--13 de la p. 78, supprimer « et $p+q=n+1$ ». Dans la
démonstration, ligne 23 de la p. 78, supprimer « et pour $n+1$ ».

\noindent\textbf{$(\mathbf{Err}_{\mathrm{III}},\,25)$} Dans (III, 1.4.15),
ligne 6 du bas de la p. 92, remplacer « \emph{de type fini} » par
« \emph{quasi-compact} ».

\noindent\textbf{$(\mathbf{Err}_{\mathrm{III}},\,26)$} Dans (III, 2.2.4),
ligne 1 de la p. 102, remplacer « \emph{que les supports de $\mathscr F$ et de
$\mathscr H$ soient propres sur $Y$} » par « \emph{que le support de
$\operatorname{Im}(\mathscr F\to\mathscr G)$ soit propre sur $Y$} ». Dans la
démonstration, remplacer le texte des lignes 10 à 16, depuis « Cela
étant\ldots{} », par :

